\subsection{Watch: Observe Sends without Replacing the Application}
\label{sec:watch}

The original \filepath{lib/watch.abclr2} is a practical debugging aid.
It wraps an evaluator, reports message-send requests, then forwards them
to the original evaluator. It does not copy the application object when
you use \code{watch}. In particular, it can observe an ordinary stateful
object without invoking the state-copying path discussed in
Section~\ref{sec:copy}.

The key distinction is \emph{outgoing sends}, not incoming requests.
Watching a caller reveals the requests it sends and the replies it sends
onward. It does not show every arrival at that caller, every operation
inside the receiver, or every host Lisp computation. The separate
\code{report} library observes a different boundary: message arrival and
script execution through a replacement metaobject.

\subsubsection{Start with One Sender}

The supplemental \filepath{tutorial/examples/12-watch.abclr2} defines
an echo target and a driver with now-, past-, and future-type sends:

\wholelisting[firstline=5,lastline=14]{tutorial/examples/12-watch.abclr2}

The file loads \code{watch} and \code{future}; the latter is needed for
this driver's future request, not for the basic watch facility.
Run this in a fresh session:

\begin{codeblock}
\begin{lstlisting}
(inc "tutorial/examples/12-watch" :use-source t)
(watch watch-driver)
[watch-driver <== [:now 7]]             ; => 7, plus trace output
(nowatch watch-driver)
[watch-driver <== [:now 8]]             ; => 8, no watch trace
\end{lstlisting}
\end{codeblock}

These are the original library's public macros. Common Lisp reader case
folding makes \code{noWatch} and \code{nowatch} the same macro name.
The trace contains a send to \code{watch-target} with value 7 and a send
of the reply onward to the driver's caller. The top-level request
\code{[:now 7]} itself was sent by the interactive evaluator, not by the
watched driver, so it is not the driver's outgoing request being reported.

The default reporter prints lines shaped like
\code{[target <== payload from sender]}. Object serial numbers are not
stable expected output. For a past-type send with a reply destination it
also prints \code{@ reply}. Its historical future display uses
\code{<<=}; that is a trace label, not a replacement for this book's
\code{[target <= payload \$ future]} source syntax.
Trace lines are asynchronous reports, not acknowledgments that a
receiver has finished processing the message.

\subsubsection{Read the Evaluator Wrapper}

Here is the complete original generator:

\wholelisting[firstline=37,lastline=55]{lib/watch.abclr2}

The dotted patterns bind the rest of the send representation as
\code{Arg}, and the remaining evaluation context as \code{Context}.
For each of the three send types, the wrapper sends the reporter four
fields: the type tag, evaluated send arguments, evaluation context, and
continuation \code{C}. It then forwards the original evaluator request
with the same continuation. The application reply path is not redirected
through the reporter.

For the standard send protocol, \code{Context} is
\code{(Env Id Gid Eval)}. Its second element is the sender's meta-level
identity; \code{[den Id]} denotes the application sender. This is how
the default formatter obtains the displayed sender. The wrapper also
answers \code{:original-evaluator}, which is useful for checking or
explicitly undoing one layer of instrumentation.

The catch-all forwards other evaluator requests unchanged. This is not
an instruction-by-instruction profiler, and operations performed by host
Lisp or a non-reifying script need not cross these evaluator patterns.
Observation adds messages and work of its own; do not infer uninstrumented
scheduling or timing from a watched run.

\subsubsection{Collect Structured Events Instead of Printed Names}

The macros select \code{default-watch-reporter}. To supply another
reporter, use the original manager's explicit protocol. The teaching
file includes this small recording object:

\wholelisting[firstline=16,lastline=22]{tutorial/examples/12-watch.abclr2}

Each event retains the same four fields the wrapper sends. The helper
\code{watch-target-events} filters those records by target identity;
it avoids comparing unstable printed names. In a fresh session after
loading the teaching file:

\begin{codeblock}
\begin{lstlisting}
[watch-manager <= [:watch-object watch-driver watch-recorder]]
[watch-driver <== [:now 7]]
[watch-driver <== [:past 8]]
[watch-driver <== [:future 9]]
[events := (watch-target-events
             [watch-recorder <== :events] watch-target)]
(mapcar #'first events)
; => (:SEND-NOW :SEND-PAST :SEND-FUTURE)
(nowatch watch-driver)
\end{lstlisting}
\end{codeblock}

There are additional events for future creation, future consumption, and
the driver's replies. Filtering for \code{watch-target} isolates the
three intended sends. The driver returns 7, \code{:submitted}, and 9:
the past-type result still acknowledges submission, not target completion.

Manager commands are past-type requests. Their handlers do not promise
application replies: do not replace \code{<=} with \code{<==} merely
to wait for installation. The bounded sessions check observable evaluator
identities and finite traces, not a general synchronization protocol.
Keep a recording reporter outside the watched sender's evaluator policy,
and bound or clear its retained events in longer experiments.

For printed output with filtering, the original \code{reporter-gen}
accepts two host Lisp functions. The predicate decides whether to print;
the formatter returns a string. For example:

\begin{codeblock}
\begin{lstlisting}
[filtered := [reporter-gen <== [:new
 :pred #'(lambda (type args context cont) (eq type :send-now))
 :format #'(lambda (type args context cont)
             (format nil "WATCH-NOW ~S" (second args)))]]]
[watch-manager <= [:watch-object watch-driver filtered]]
[watch-driver <== [:now :probe]]
(nowatch watch-driver)
\end{lstlisting}
\end{codeblock}

The formatter sees evaluated arguments, not source text. These functions
are host lambdas: do not place ABCL/R2 sends in their bodies or assume they
can capture arbitrary surrounding ABCL bindings. Use ordinary Lisp
parameters or explicit Lisp closure construction as in the earlier
lambda-boundary discussion.

\subsubsection{Choose the Scope and Understand Restoration}

\begin{longtable}{L{0.33\textwidth}L{0.57\textwidth}}
\textbf{Public operation} & \textbf{Scope of the change}\\\hline
\code{(watch obj)} / \code{(nowatch obj)} & Replaces/restores this object's
evaluator through its metaobject. Requires compatible reflective operations.\\
\code{(watch-group g)} / \code{(nowatch-group g)} & Changes the group's
evaluator for subsequent object creation; does not rewrite existing objects.\\
\code{(watch-toplevel)} / \code{(nowatch-toplevel)} & Changes/restores
\code{*eval*}, the interactive evaluator; not every object's evaluator.\\
\end{longtable}

The manager saves the previous evaluator before installing the wrapper.
For an object it sends \code{[:set-evaluator wrapper]} to the metaobject;
for a group it sends \code{[:replace-evaluator wrapper]} to the group
manager. The top-level operation assigns \code{*eval*}. These are three
different places to install the same kind of wrapper.

In a group experiment, create one sender before \code{watch-group} and
another after it. The first retains its original evaluator; the second
captures the wrapper. After \code{nowatch-group}, a third sender gets the
restored group policy, but the second still holds its wrapper. Group
restoration is not a broadcast repair of all existing objects.

There is also an original, preconfigured \code{watching} group. An object
defined with \code{(group watching)} uses its tracing evaluator from
creation. Alternatively, after loading \code{copy}, the library suggests
\code{(copy name :group watching)}. That rebinds \code{name} to a
\emph{different object}; it is not equivalent to \code{(watch name)}.
All the queue-sharing and local-state-copy restrictions of
Section~\ref{sec:copy} still apply. The checked copy example has no local
state entries and does not establish general copying compatibility.

Wrappers obtained from a group's creation policy are not entries in the
manager's watched-object list. Thus \code{(nowatch obj)} alone cannot
remove such a preconfigured wrapper. For a known single wrapper, an
explicit object-level restoration is:

\begin{codeblock}
\begin{lstlisting}
[wrapper := [[meta obj] <== :evaluator]]
[underlying := [wrapper <== :original-evaluator]]
[[meta obj] <= [:set-evaluator underlying]]
\end{lstlisting}
\end{codeblock}

Use this only when you know that the current evaluator implements this
watch-wrapper protocol. It removes one layer from that object, not the
group policy. Do not replace an unrelated custom evaluator blindly.

\paragraph{Avoid repeated installation.}
The original manager does not make \code{watch} idempotent. Calling it
twice adds two wrappers. On disabling the object, it restores the most
recent saved evaluator but deletes all saved entries for that object.
A second \code{nowatch} therefore leaves the inner wrapper installed.
The tests observe this boundary and explicitly restore the evaluator
saved before the experiment. Pair one installation with one restoration;
do not rely on repeated object, group, or top-level requests behaving as
a balanced instrumentation stack.

\paragraph{Evidence and exercise.}
Run \code{make -C tutorial watch-check}. Forty-one probes check
now/past/future results and event payloads, sender context, default and
filtered reporters, restoration, group creation policy, top-level control,
the preconfigured group, stateful watching without copying, an
empty-local-state watched copy, and repeated
installation boundaries. Libraries are copied unchanged; transcripts and
generated files stay in \filepath{build/watch-tests/}.

Why watch the driver rather than just the echo target? The driver issues
the application requests of interest. Why does turning off the group not
silence every existing sender? Each already has an evaluator reference.
As an exercise, record only \code{:send-past} events and explain which
are application requests and which are replies forwarded by the driver.
